\section{四年学习路线}
\label{sec:01-system-overview-03}

四年路线按"同一条主线逐层加码"组织：每一年都在同一个系统上增加一层能力，
而不是换一个互不相关的题目。示意安排见\cref{tab:roadmap}。

\begin{table}[htbp]
\centering
\caption{四年学习路线示意}
\label{tab:roadmap}
\begin{tabularx}{\linewidth}{@{}P{14mm}P{40mm}Y@{}}
\toprule
学年 & 主线 & 可验收的产物 \\
\midrule
大一 & 数字逻辑与计算机组成 & 用硬件描述语言实现单周期处理器，仿真跑通一段汇编 \\
大二 & 指令集与微架构 & 五级流水线上板运行，正确处理冒险与旁路 \\
大三 & 工具链与操作系统 & 在自研处理器核上启动实时操作系统或精简 Linux \\
大四 & SoC 与智能应用 & 在自研加速器上跑通量化模型的 TinyML 推理 \\
\bottomrule
\end{tabularx}
\end{table}

% TODO：撰写本节正文。建议结构：
%   1) 问题引入与学习目标
%   2) 原理与机制（配图、公式）
%   3) 示例或实验步骤
%   4) 小结与思考题
本节待撰写。
